Em 6 de outubro de 2026, o teste de equivalência enviou um conjunto determinístico de 50 registros a cada variante, em execuções separadas. Os UUIDs, campos, instantes de ocorrência e payloads foram iguais; somente o identificador de execução, transportado fora do conteúdo canônico, diferiu. Ambas as exportações continham 50 linhas. A comparação independente confirmou igualdade dos identificadores, campos, payloads e SHA-256, desconsiderando apenas \texttt{persisted\_at}, pois as inserções ocorreram em momentos distintos.

O reenvio do primeiro registro, com o mesmo conteúdo, produziu um marco de persistência \texttt{duplicate} em cada variante e não aumentou a contagem de linhas. Em seguida, o payload desse UUID foi alterado de zero para nove. A API síncrona respondeu 409. A assíncrona respondeu 202 por confirmar a publicação; o consumidor identificou \texttt{event\_id\_content\_conflict} e confirmou o encaminhamento da mensagem conflitante à DLQ. O conteúdo original no banco permaneceu inalterado.

Esse resultado evidencia equivalência funcional para o conjunto examinado e verifica os caminhos de idempotência e conflito. Não demonstra equivalência para todos os possíveis payloads ou condições de concorrência e não constitui comparação de desempenho. As execuções válidas foram identificadas pelos prefixos \texttt{e71d7147} e \texttt{6cdf03fd}; os UUIDs completos, as exportações e seus checksums constam do catálogo técnico de evidências.
